\documentclass[a4paper]{article}

\def\nbook {Introduction to Quantum Mechanics \\ Second Edition}
\def\nbookshort {Griffiths Quantum Mechanics}

% Shared preamble. The fallback lets the document build whether the working
% directory is its own folder or the project root.
\IfFileExists{headerintroquantummechanics.tex}{%%%%%%%%%%%%%%%%%%%%
% AUTHOR AND HEADER
% define \nbook
%%%%%%%%%%%%%%%%%%%%

\makeatletter
\ifx \nauthor\undefined
	\def\nauthor{David A. Lee}
\else
\fi

\author{\nbook \\ \small David J. Griffiths \\ \small Solutions by \nauthor}
\date{}

%%%%%%%%%%%%%%%%%%%%
% PACKAGES
%%%%%%%%%%%%%%%%%%%%

\usepackage{alltt} 
\usepackage{amsfonts}
\usepackage{amsmath}
\usepackage{amssymb}
\usepackage{amsthm}
\usepackage{bm}
\usepackage{booktabs}
\usepackage{braket} % \bra, \ket, \braket; not `physics`, which clashes with \abs/\norm below
\usepackage{caption}
\usepackage{centernot}
\usepackage{enumitem}
\usepackage{fancyhdr}
\usepackage[T1]{fontenc}
\usepackage{graphicx}
% The documents sit one level below the project root, where the chapN/ figure
% folders live. Listing both "./" and "../"
% lets those be referenced the same way regardless of where the build runs.
\graphicspath{{./}{../}}
\usepackage{mathdots}
\usepackage{mathtools}
% \usepackage{microtype}
\usepackage{multirow}
\usepackage{pdflscape}
\usepackage{pgfplots}
\pgfplotsset{compat=1.18}
\usepackage{siunitx}
\usepackage{slashed}
\usepackage{tabularx}
\usepackage{tikz}
\usepackage{titletoc}
\usepackage{tkz-euclide}
\usepackage[normalem]{ulem}
\usepackage{url}
\usepackage[all]{xy}
\usepackage{imakeidx}

% fonts
% use \textsf for the sans family (problem titles, section headings)
% Latin Modern Sans, not CM Bright: cmbr under T1 has no Type 1 outlines, so it
% renders as Type 3 bitmaps, which TeXpresso's viewer cannot draw at all.
\renewcommand{\sfdefault}{lmss} % keeps default font Computer Modern Roman

% makeindex for table of contents
\makeindex[intoc, title=Index]
\indexsetup{othercode={\lhead{ Index }}}

% table of contents remove numbering
\usepackage{tocloft} % for coloring subsections -- note that sections have standard document hyperlink coloring, i.e. color = doc
\setcounter{secnumdepth}{0} % remove numbering
\renewcommand{\cftsubsecfont}{\hypersetup{linkcolor=subsect!90}} % subsection coloring
\renewcommand{\cftsubsubsecfont}{\hypersetup{linkcolor=subsubsect!90}} % subsubsection coloring
\renewcommand{\contentsname}{\textsf{Table of Contents}}

% pagestyle
\pagestyle{fancyplain}

% header
\lhead{	\nouppercase{\leftmark} }
\ifx \nextra \undefined
  \rhead{
    \ifnum\thepage=1
    \else
      \nbookshort \  | \nauthor
    \fi}
\else
  \rhead{
    \ifnum\thepage=1
    \else
      \nbookshort (\nextra)
    \fi}
\fi

% proof environment
\newcommand{\prooffont}{\scshape}
\usepackage{xpatch}
\tracingpatches
\xpatchcmd{\proof}{\itshape}{\prooffont}{}{}

% Python environment (for typesetting Python code)

% Default fixed font does not support bold face
\DeclareFixedFont{\ttb}{T1}{txtt}{bx}{n}{8} % for bold
\DeclareFixedFont{\ttm}{T1}{txtt}{m}{n}{8}  % for normal

% Custom colors
\usepackage{color}
\definecolor{deepblue}{rgb}{0,0,0.5}
\definecolor{deepred}{rgb}{0.6,0,0}
\definecolor{deepgreen}{rgb}{0,0.5,0}

\usepackage{listings}

% Python style for highlighting
\newcommand\pythonstyle{\lstset{
language=Python,
basicstyle=\ttm,
morekeywords={self},              % Add keywords here
keywordstyle=\ttb\color{deepblue},
emph={MyClass,__init__},          % Custom highlighting
emphstyle=\ttb\color{deepred},    % Custom highlighting style
stringstyle=\color{deepgreen},
frame=tb,                         % Any extra options here
showstringspaces=false
}}


% Python environment
\lstnewenvironment{python}[1][]
{
\pythonstyle
\lstset{#1}
}
{}

% Python for external files
\newcommand\pythonexternal[2][]{{
\pythonstyle
\lstinputlisting[#1]{#2}}}

% Python for inline
\newcommand\pythoninline[1]{{\pythonstyle\lstinline!#1!}}

%%%%%%%%%%%%%
% FUNCTIONS %
%%%%%%%%%%%%%

\DeclarePairedDelimiter\abs{\lvert}{\rvert}%
% plain \abs{x} autosizes like \left| x \right|; \abs*{x} keeps fixed-size bars
\DeclarePairedDelimiter\absfixed{\lvert}{\rvert}%
\RenewDocumentCommand{\abs}{s}{\IfBooleanTF{#1}{\absfixed}{\absfixed*}}%
\DeclarePairedDelimiter\norm{\lVert}{\rVert}%
% plain \bra, \ket, \braket autosize like braket's \Bra, \Ket, \Braket (including
% the | separators); starred forms keep braket's fixed-size delimiters
\let\brafixed\bra
\let\ketfixed\ket
\let\braketfixed\braket
\RenewDocumentCommand{\bra}{s m}{\IfBooleanTF{#1}{\brafixed{#2}}{\Bra{#2}}}%
\RenewDocumentCommand{\ket}{s m}{\IfBooleanTF{#1}{\ketfixed{#2}}{\Ket{#2}}}%
\RenewDocumentCommand{\braket}{s m}{\IfBooleanTF{#1}{\braketfixed{#2}}{\Braket{#2}}}%

%%%%%%%%%%%%%%%%%%%%
% DOCUMENT GEOMETRY
%%%%%%%%%%%%%%%%%%%%

\ifx \ntrim \undefined
\else
  \usepackage{geometry}
  \geometry{
	  papersize={379pt, 699pt},
	  textwidth=345pt,
	  left=17pt,
	  top=54pt,
	  right=17pt,
  }
\fi

% maketitle statement
\title{}
\ifx \nisofficial \undefined
\let\@real@maketitle\maketitle
\renewcommand{\maketitle}{\@real@maketitle\begin{center}\begin{minipage}[c]{0.9\textwidth}\centering\footnotesize All errors, typographical and substantive, and other offenses, are entirely my own.\end{minipage}\end{center}}
\else
\fi

%%%%%%%%%%%%%%%%%%%%
% THEOREMS
%%%%%%%%%%%%%%%%%%%%

%\theoremstyle{definition}
\newtheorem*{axiom}{Axiom}
\newtheorem*{claim}{Claim}
\newtheorem*{conjecture}{Conjecture}
\newtheorem*{cor}{Corollary}
%\newtheorem*{dfn}{Definition}
\newtheorem*{eg}{Example}
\newtheorem*{ex}{Exercise}
%\newtheorem*{lemma}{Lemma}
\newtheorem*{prop}{Proposition}
%\newtheorem*{thm}{Theorem}
\newtheorem*{remark}{Remark}
\newtheorem*{warning}{Warning}

%%%%%%%%%%%%%%%%%%%%
% FORMATTING AESTHETICS
%%%%%%%%%%%%%%%%%%%%

% itemize bullets
\renewcommand{\labelitemi}{-}
\renewcommand{\labelitemii}{$\circ$}
\renewcommand{\labelenumi}{(\roman{*})}

% new page section
%\let\stdsection\section
%\renewcommand\section{\newpage\stdsection}

% new page subsection
%\let\stdsubsection\subsection
%\renewcommand\subsection{\newpage\stdsubsection}

%%%%%%%%%%%%%%%%%%%%
% \mathbb commands
%%%%%%%%%%%%%%%%%%%%

\newcommand{\C}{\mathbb{C}}
\newcommand{\N}{\mathbb{N}}
\newcommand{\Q}{\mathbb{Q}}
\newcommand{\R}{\mathbb{R}}
\newcommand{\Z}{\mathbb{Z}}

%%%%%%%%%%%%%%%%%%%%
% Complex Numbers
%%%%%%%%%%%%%%%%%%%%

\DeclareMathOperator{\Real}{Re}
\DeclareMathOperator{\Imag}{Im}

%%%%%%%%%%%%%%%%%%%%
% Optimization
%%%%%%%%%%%%%%%%%%%%

\DeclareMathOperator*{\argmax}{arg\,max}

%%%%%%%%%%%%%%%%%%%%
% Probability Operators
%%%%%%%%%%%%%%%%%%%%

\DeclareMathOperator{\Bernoulli}{Bernoulli}
\DeclareMathOperator{\betaD}{Beta}
\DeclareMathOperator{\bias}{\textsf{bias}}
\DeclareMathOperator{\binomial}{Binomial}
\DeclareMathOperator{\corr}{Corr}
\DeclareMathOperator{\cov}{\textsf{Cov}}
\DeclareMathOperator{\Exp}{Exp}
\DeclareMathOperator{\gammaD}{Gamma}
\DeclareMathOperator{\mse}{\textsf{MSE}}
\DeclareMathOperator{\multinomial}{Multinomial}
\DeclareMathOperator{\Poisson}{Poisson}
\DeclareMathOperator{\sd}{sd}
\DeclareMathOperator{\se}{\textsf{se}}
\DeclareMathOperator{\Uniform}{Uniform}
\newcommand{\Prob}{\mathbb{P}}
\newcommand{\var}{\mathbb{V}}
\newcommand{\E}{\mathbb{E}}

%%%%%%%%%%%%%%%%%%%%
% TCOLORBOX CODE FROM SeniorMars
%%%%%%%%%%%%%%%%%%%%

%\usepackage[varbb]{newpxmath} changes math font to newpxmath
\usepackage{xfrac}
\usepackage[makeroom]{cancel}
\usepackage{mathtools}
\usepackage{bookmark}
\usepackage{enumitem}
\usepackage{hyperref}
\usepackage{theoremref}
\hypersetup{
	pdftitle={Assignment},
	colorlinks=true, linkcolor=doc!90,
	bookmarksnumbered=true,
	bookmarksopen=true
}
\usepackage[most,many,breakable]{tcolorbox}
\usepackage{xcolor}
\usepackage{varwidth}
\usepackage{varwidth}
\usepackage{etoolbox}
%\usepackage{authblk}
\usepackage{nameref}
\usepackage{multicol,array}
\usepackage{tikz-cd}
\usepackage[ruled,vlined,linesnumbered]{algorithm2e}
\usepackage{comment} % enables the use of multi-line comments (\ifx \fi)
\usepackage{import}
\usepackage{xifthen}
\usepackage{pdfpages}
\usepackage{transparent}

\newcommand\mycommfont[1]{\footnotesize\ttfamily\textcolor{blue}{#1}}
\SetCommentSty{mycommfont}
\newcommand{\incfig}[1]{%
    \def\svgwidth{\columnwidth}
    \import{./figures/}{#1.pdf_tex}
}

\usepackage{tikzsymbols}

% colors -- change them as you may!

\definecolor{myg}{RGB}{56, 140, 70}
\definecolor{myb}{RGB}{45, 111, 177}
\definecolor{myr}{RGB}{199, 68, 64}
\definecolor{mytheorembg}{HTML}{F2F2F9}
\definecolor{mytheoremfr}{HTML}{00007B}
\definecolor{mylenmabg}{HTML}{FFFAF8}
\definecolor{mylenmafr}{HTML}{983b0f}
\definecolor{mypropbg}{HTML}{f2fbfc}
\definecolor{mypropfr}{HTML}{191971}
\definecolor{myexamplebg}{HTML}{F2FBF8}
\definecolor{myexamplefr}{HTML}{88D6D1}
\definecolor{myexampleti}{HTML}{2A7F7F}
\definecolor{mydefinitbg}{HTML}{E5E5FF}
\definecolor{mydefinitfr}{HTML}{3F3FA3}
\definecolor{notesgreen}{RGB}{0,162,0}
\definecolor{myp}{RGB}{197, 92, 212}
\definecolor{mygr}{HTML}{880808}
\definecolor{myred}{RGB}{127,0,0}
\definecolor{myyellow}{RGB}{169,121,69}
\definecolor{myexercisebg}{HTML}{F2FBF8}
\definecolor{myexercisefg}{HTML}{88D6D1}
\definecolor{doc}{RGB}{0,0,255}
\definecolor{subsect}{RGB}{255,0,0}
\definecolor{subsubsect}{RGB}{128,0,255}
\definecolor{question}{HTML}{024593}

%================================
% THEOREM BOX
%================================

%\providecommand\theoremnumber{}
\tcbuselibrary{theorems,skins,hooks}
\newtcbtheorem{Theorem}{Theorem}
{%
	enhanced,
	breakable,
	colback = mytheorembg,
	frame hidden,
	boxrule = 0sp,
	borderline west = {2pt}{0pt}{mytheoremfr},
	sharp corners,
	detach title,
	before upper = \tcbtitle\par\smallskip,
	coltitle = mytheoremfr,
	fonttitle = \bfseries\sffamily,
	description font = \mdseries,
	separator sign none,
	segmentation style={solid, mytheoremfr},
}
{th}


%================================
% Question BOX
%================================

\makeatletter
\newtcbtheorem{question}{Problem}{enhanced,
	breakable,
	colback=white,
	colframe=question,
	attach boxed title to top left={yshift*=-\tcboxedtitleheight},
	fonttitle=\bfseries\sffamily,
	title={#2},
	boxed title size=title,
	boxed title style={%
			sharp corners,
			rounded corners=northwest,
			colback=tcbcolframe,
			boxrule=0pt,
		},
	underlay boxed title={%
			\path[fill=tcbcolframe] (title.south west)--(title.south east)
			to[out=0, in=180] ([xshift=5mm]title.east)--
			(title.center-|frame.east)
			[rounded corners=\kvtcb@arc] |-
			(frame.north) -| cycle;
		},
	#1
}{def}
\makeatother

%================================
% PROPERTIES BOX
%================================

\tcbuselibrary{theorems,skins,hooks}
\newtcbtheorem{Property}{Property}
{%
	enhanced,
	breakable,
	colback = mytheorembg,
	frame hidden,
	boxrule = 0sp,
	borderline west = {2pt}{0pt}{mytheoremfr},
	sharp corners,
	detach title,
	before upper = \tcbtitle\par\smallskip,
	coltitle = mytheoremfr,
	fonttitle = \bfseries\sffamily,
	description font = \mdseries,
	separator sign none,
	segmentation style={solid, mytheoremfr},
}
{th}

%================================
% LEMMA
%================================

\tcbuselibrary{theorems,skins,hooks}
\newtcbtheorem[number within=section]{Lemma}{Lemma}
{%
	enhanced,
	breakable,
	colback = mylenmabg,
	frame hidden,
	boxrule = 0sp,
	borderline west = {2pt}{0pt}{mylenmafr},
	sharp corners,
	detach title,
	before upper = \tcbtitle\par\smallskip,
	coltitle = mylenmafr,
	fonttitle = \bfseries\sffamily,
	description font = \mdseries,
	separator sign none,
	segmentation style={solid, mylenmafr},
}
{th}


% Commands for special boxes
\newcommand{\thm}[2]{\begin{Theorem*}{#1}{}#2\end{Theorem*}}
\newcommand{\qs}[2]{\begin{question*}{#1}{}#2\end{question*}}
\newcommand{\prp}[2]{\begin{Property*}{#1}{}#2\end{Property*}}
\newcommand{\lemma}[2]{\begin{Lemma*}{#1}{}#2\end{Lemma*}}


%%%%%%%%%%%%%%%%%%%%
% SHARED DOCUMENT SETTINGS
%%%%%%%%%%%%%%%%%%%%
% Moved here from the chapter document's preamble so that every document
% (chapters and the merged manual) picks up the same settings from one place.

\hfuzz=125pt      % ignore overfull boxes up to 5pt
\hbadness=10000   % suppress underfull warnings

\usetikzlibrary{arrows.meta, calc, decorations.markings, backgrounds}
\usepackage{tikz-3dplot}
\usepackage{float}
\usepackage{esint}
}{%%%%%%%%%%%%%%%%%%%%
% AUTHOR AND HEADER
% define \nbook
%%%%%%%%%%%%%%%%%%%%

\makeatletter
\ifx \nauthor\undefined
	\def\nauthor{David A. Lee}
\else
\fi

\author{\nbook \\ \small David J. Griffiths \\ \small Solutions by \nauthor}
\date{}

%%%%%%%%%%%%%%%%%%%%
% PACKAGES
%%%%%%%%%%%%%%%%%%%%

\usepackage{alltt} 
\usepackage{amsfonts}
\usepackage{amsmath}
\usepackage{amssymb}
\usepackage{amsthm}
\usepackage{bm}
\usepackage{booktabs}
\usepackage{braket} % \bra, \ket, \braket; not `physics`, which clashes with \abs/\norm below
\usepackage{caption}
\usepackage{centernot}
\usepackage{enumitem}
\usepackage{fancyhdr}
\usepackage[T1]{fontenc}
\usepackage{graphicx}
% The documents sit one level below the project root, where the chapN/ figure
% folders live. Listing both "./" and "../"
% lets those be referenced the same way regardless of where the build runs.
\graphicspath{{./}{../}}
\usepackage{mathdots}
\usepackage{mathtools}
% \usepackage{microtype}
\usepackage{multirow}
\usepackage{pdflscape}
\usepackage{pgfplots}
\pgfplotsset{compat=1.18}
\usepackage{siunitx}
\usepackage{slashed}
\usepackage{tabularx}
\usepackage{tikz}
\usepackage{titletoc}
\usepackage{tkz-euclide}
\usepackage[normalem]{ulem}
\usepackage{url}
\usepackage[all]{xy}
\usepackage{imakeidx}

% fonts
% use \textsf for the sans family (problem titles, section headings)
% Latin Modern Sans, not CM Bright: cmbr under T1 has no Type 1 outlines, so it
% renders as Type 3 bitmaps, which TeXpresso's viewer cannot draw at all.
\renewcommand{\sfdefault}{lmss} % keeps default font Computer Modern Roman

% makeindex for table of contents
\makeindex[intoc, title=Index]
\indexsetup{othercode={\lhead{ Index }}}

% table of contents remove numbering
\usepackage{tocloft} % for coloring subsections -- note that sections have standard document hyperlink coloring, i.e. color = doc
\setcounter{secnumdepth}{0} % remove numbering
\renewcommand{\cftsubsecfont}{\hypersetup{linkcolor=subsect!90}} % subsection coloring
\renewcommand{\cftsubsubsecfont}{\hypersetup{linkcolor=subsubsect!90}} % subsubsection coloring
\renewcommand{\contentsname}{\textsf{Table of Contents}}

% pagestyle
\pagestyle{fancyplain}

% header
\lhead{	\nouppercase{\leftmark} }
\ifx \nextra \undefined
  \rhead{
    \ifnum\thepage=1
    \else
      \nbookshort \  | \nauthor
    \fi}
\else
  \rhead{
    \ifnum\thepage=1
    \else
      \nbookshort (\nextra)
    \fi}
\fi

% proof environment
\newcommand{\prooffont}{\scshape}
\usepackage{xpatch}
\tracingpatches
\xpatchcmd{\proof}{\itshape}{\prooffont}{}{}

% Python environment (for typesetting Python code)

% Default fixed font does not support bold face
\DeclareFixedFont{\ttb}{T1}{txtt}{bx}{n}{8} % for bold
\DeclareFixedFont{\ttm}{T1}{txtt}{m}{n}{8}  % for normal

% Custom colors
\usepackage{color}
\definecolor{deepblue}{rgb}{0,0,0.5}
\definecolor{deepred}{rgb}{0.6,0,0}
\definecolor{deepgreen}{rgb}{0,0.5,0}

\usepackage{listings}

% Python style for highlighting
\newcommand\pythonstyle{\lstset{
language=Python,
basicstyle=\ttm,
morekeywords={self},              % Add keywords here
keywordstyle=\ttb\color{deepblue},
emph={MyClass,__init__},          % Custom highlighting
emphstyle=\ttb\color{deepred},    % Custom highlighting style
stringstyle=\color{deepgreen},
frame=tb,                         % Any extra options here
showstringspaces=false
}}


% Python environment
\lstnewenvironment{python}[1][]
{
\pythonstyle
\lstset{#1}
}
{}

% Python for external files
\newcommand\pythonexternal[2][]{{
\pythonstyle
\lstinputlisting[#1]{#2}}}

% Python for inline
\newcommand\pythoninline[1]{{\pythonstyle\lstinline!#1!}}

%%%%%%%%%%%%%
% FUNCTIONS %
%%%%%%%%%%%%%

\DeclarePairedDelimiter\abs{\lvert}{\rvert}%
% plain \abs{x} autosizes like \left| x \right|; \abs*{x} keeps fixed-size bars
\DeclarePairedDelimiter\absfixed{\lvert}{\rvert}%
\RenewDocumentCommand{\abs}{s}{\IfBooleanTF{#1}{\absfixed}{\absfixed*}}%
\DeclarePairedDelimiter\norm{\lVert}{\rVert}%
% plain \bra, \ket, \braket autosize like braket's \Bra, \Ket, \Braket (including
% the | separators); starred forms keep braket's fixed-size delimiters
\let\brafixed\bra
\let\ketfixed\ket
\let\braketfixed\braket
\RenewDocumentCommand{\bra}{s m}{\IfBooleanTF{#1}{\brafixed{#2}}{\Bra{#2}}}%
\RenewDocumentCommand{\ket}{s m}{\IfBooleanTF{#1}{\ketfixed{#2}}{\Ket{#2}}}%
\RenewDocumentCommand{\braket}{s m}{\IfBooleanTF{#1}{\braketfixed{#2}}{\Braket{#2}}}%

%%%%%%%%%%%%%%%%%%%%
% DOCUMENT GEOMETRY
%%%%%%%%%%%%%%%%%%%%

\ifx \ntrim \undefined
\else
  \usepackage{geometry}
  \geometry{
	  papersize={379pt, 699pt},
	  textwidth=345pt,
	  left=17pt,
	  top=54pt,
	  right=17pt,
  }
\fi

% maketitle statement
\title{}
\ifx \nisofficial \undefined
\let\@real@maketitle\maketitle
\renewcommand{\maketitle}{\@real@maketitle\begin{center}\begin{minipage}[c]{0.9\textwidth}\centering\footnotesize All errors, typographical and substantive, and other offenses, are entirely my own.\end{minipage}\end{center}}
\else
\fi

%%%%%%%%%%%%%%%%%%%%
% THEOREMS
%%%%%%%%%%%%%%%%%%%%

%\theoremstyle{definition}
\newtheorem*{axiom}{Axiom}
\newtheorem*{claim}{Claim}
\newtheorem*{conjecture}{Conjecture}
\newtheorem*{cor}{Corollary}
%\newtheorem*{dfn}{Definition}
\newtheorem*{eg}{Example}
\newtheorem*{ex}{Exercise}
%\newtheorem*{lemma}{Lemma}
\newtheorem*{prop}{Proposition}
%\newtheorem*{thm}{Theorem}
\newtheorem*{remark}{Remark}
\newtheorem*{warning}{Warning}

%%%%%%%%%%%%%%%%%%%%
% FORMATTING AESTHETICS
%%%%%%%%%%%%%%%%%%%%

% itemize bullets
\renewcommand{\labelitemi}{-}
\renewcommand{\labelitemii}{$\circ$}
\renewcommand{\labelenumi}{(\roman{*})}

% new page section
%\let\stdsection\section
%\renewcommand\section{\newpage\stdsection}

% new page subsection
%\let\stdsubsection\subsection
%\renewcommand\subsection{\newpage\stdsubsection}

%%%%%%%%%%%%%%%%%%%%
% \mathbb commands
%%%%%%%%%%%%%%%%%%%%

\newcommand{\C}{\mathbb{C}}
\newcommand{\N}{\mathbb{N}}
\newcommand{\Q}{\mathbb{Q}}
\newcommand{\R}{\mathbb{R}}
\newcommand{\Z}{\mathbb{Z}}

%%%%%%%%%%%%%%%%%%%%
% Complex Numbers
%%%%%%%%%%%%%%%%%%%%

\DeclareMathOperator{\Real}{Re}
\DeclareMathOperator{\Imag}{Im}

%%%%%%%%%%%%%%%%%%%%
% Optimization
%%%%%%%%%%%%%%%%%%%%

\DeclareMathOperator*{\argmax}{arg\,max}

%%%%%%%%%%%%%%%%%%%%
% Probability Operators
%%%%%%%%%%%%%%%%%%%%

\DeclareMathOperator{\Bernoulli}{Bernoulli}
\DeclareMathOperator{\betaD}{Beta}
\DeclareMathOperator{\bias}{\textsf{bias}}
\DeclareMathOperator{\binomial}{Binomial}
\DeclareMathOperator{\corr}{Corr}
\DeclareMathOperator{\cov}{\textsf{Cov}}
\DeclareMathOperator{\Exp}{Exp}
\DeclareMathOperator{\gammaD}{Gamma}
\DeclareMathOperator{\mse}{\textsf{MSE}}
\DeclareMathOperator{\multinomial}{Multinomial}
\DeclareMathOperator{\Poisson}{Poisson}
\DeclareMathOperator{\sd}{sd}
\DeclareMathOperator{\se}{\textsf{se}}
\DeclareMathOperator{\Uniform}{Uniform}
\newcommand{\Prob}{\mathbb{P}}
\newcommand{\var}{\mathbb{V}}
\newcommand{\E}{\mathbb{E}}

%%%%%%%%%%%%%%%%%%%%
% TCOLORBOX CODE FROM SeniorMars
%%%%%%%%%%%%%%%%%%%%

%\usepackage[varbb]{newpxmath} changes math font to newpxmath
\usepackage{xfrac}
\usepackage[makeroom]{cancel}
\usepackage{mathtools}
\usepackage{bookmark}
\usepackage{enumitem}
\usepackage{hyperref}
\usepackage{theoremref}
\hypersetup{
	pdftitle={Assignment},
	colorlinks=true, linkcolor=doc!90,
	bookmarksnumbered=true,
	bookmarksopen=true
}
\usepackage[most,many,breakable]{tcolorbox}
\usepackage{xcolor}
\usepackage{varwidth}
\usepackage{varwidth}
\usepackage{etoolbox}
%\usepackage{authblk}
\usepackage{nameref}
\usepackage{multicol,array}
\usepackage{tikz-cd}
\usepackage[ruled,vlined,linesnumbered]{algorithm2e}
\usepackage{comment} % enables the use of multi-line comments (\ifx \fi)
\usepackage{import}
\usepackage{xifthen}
\usepackage{pdfpages}
\usepackage{transparent}

\newcommand\mycommfont[1]{\footnotesize\ttfamily\textcolor{blue}{#1}}
\SetCommentSty{mycommfont}
\newcommand{\incfig}[1]{%
    \def\svgwidth{\columnwidth}
    \import{./figures/}{#1.pdf_tex}
}

\usepackage{tikzsymbols}

% colors -- change them as you may!

\definecolor{myg}{RGB}{56, 140, 70}
\definecolor{myb}{RGB}{45, 111, 177}
\definecolor{myr}{RGB}{199, 68, 64}
\definecolor{mytheorembg}{HTML}{F2F2F9}
\definecolor{mytheoremfr}{HTML}{00007B}
\definecolor{mylenmabg}{HTML}{FFFAF8}
\definecolor{mylenmafr}{HTML}{983b0f}
\definecolor{mypropbg}{HTML}{f2fbfc}
\definecolor{mypropfr}{HTML}{191971}
\definecolor{myexamplebg}{HTML}{F2FBF8}
\definecolor{myexamplefr}{HTML}{88D6D1}
\definecolor{myexampleti}{HTML}{2A7F7F}
\definecolor{mydefinitbg}{HTML}{E5E5FF}
\definecolor{mydefinitfr}{HTML}{3F3FA3}
\definecolor{notesgreen}{RGB}{0,162,0}
\definecolor{myp}{RGB}{197, 92, 212}
\definecolor{mygr}{HTML}{880808}
\definecolor{myred}{RGB}{127,0,0}
\definecolor{myyellow}{RGB}{169,121,69}
\definecolor{myexercisebg}{HTML}{F2FBF8}
\definecolor{myexercisefg}{HTML}{88D6D1}
\definecolor{doc}{RGB}{0,0,255}
\definecolor{subsect}{RGB}{255,0,0}
\definecolor{subsubsect}{RGB}{128,0,255}
\definecolor{question}{HTML}{024593}

%================================
% THEOREM BOX
%================================

%\providecommand\theoremnumber{}
\tcbuselibrary{theorems,skins,hooks}
\newtcbtheorem{Theorem}{Theorem}
{%
	enhanced,
	breakable,
	colback = mytheorembg,
	frame hidden,
	boxrule = 0sp,
	borderline west = {2pt}{0pt}{mytheoremfr},
	sharp corners,
	detach title,
	before upper = \tcbtitle\par\smallskip,
	coltitle = mytheoremfr,
	fonttitle = \bfseries\sffamily,
	description font = \mdseries,
	separator sign none,
	segmentation style={solid, mytheoremfr},
}
{th}


%================================
% Question BOX
%================================

\makeatletter
\newtcbtheorem{question}{Problem}{enhanced,
	breakable,
	colback=white,
	colframe=question,
	attach boxed title to top left={yshift*=-\tcboxedtitleheight},
	fonttitle=\bfseries\sffamily,
	title={#2},
	boxed title size=title,
	boxed title style={%
			sharp corners,
			rounded corners=northwest,
			colback=tcbcolframe,
			boxrule=0pt,
		},
	underlay boxed title={%
			\path[fill=tcbcolframe] (title.south west)--(title.south east)
			to[out=0, in=180] ([xshift=5mm]title.east)--
			(title.center-|frame.east)
			[rounded corners=\kvtcb@arc] |-
			(frame.north) -| cycle;
		},
	#1
}{def}
\makeatother

%================================
% PROPERTIES BOX
%================================

\tcbuselibrary{theorems,skins,hooks}
\newtcbtheorem{Property}{Property}
{%
	enhanced,
	breakable,
	colback = mytheorembg,
	frame hidden,
	boxrule = 0sp,
	borderline west = {2pt}{0pt}{mytheoremfr},
	sharp corners,
	detach title,
	before upper = \tcbtitle\par\smallskip,
	coltitle = mytheoremfr,
	fonttitle = \bfseries\sffamily,
	description font = \mdseries,
	separator sign none,
	segmentation style={solid, mytheoremfr},
}
{th}

%================================
% LEMMA
%================================

\tcbuselibrary{theorems,skins,hooks}
\newtcbtheorem[number within=section]{Lemma}{Lemma}
{%
	enhanced,
	breakable,
	colback = mylenmabg,
	frame hidden,
	boxrule = 0sp,
	borderline west = {2pt}{0pt}{mylenmafr},
	sharp corners,
	detach title,
	before upper = \tcbtitle\par\smallskip,
	coltitle = mylenmafr,
	fonttitle = \bfseries\sffamily,
	description font = \mdseries,
	separator sign none,
	segmentation style={solid, mylenmafr},
}
{th}


% Commands for special boxes
\newcommand{\thm}[2]{\begin{Theorem*}{#1}{}#2\end{Theorem*}}
\newcommand{\qs}[2]{\begin{question*}{#1}{}#2\end{question*}}
\newcommand{\prp}[2]{\begin{Property*}{#1}{}#2\end{Property*}}
\newcommand{\lemma}[2]{\begin{Lemma*}{#1}{}#2\end{Lemma*}}


%%%%%%%%%%%%%%%%%%%%
% SHARED DOCUMENT SETTINGS
%%%%%%%%%%%%%%%%%%%%
% Moved here from the chapter document's preamble so that every document
% (chapters and the merged manual) picks up the same settings from one place.

\hfuzz=125pt      % ignore overfull boxes up to 5pt
\hbadness=10000   % suppress underfull warnings

\usetikzlibrary{arrows.meta, calc, decorations.markings, backgrounds}
\usepackage{tikz-3dplot}
\usepackage{float}
\usepackage{esint}
}

\begin{document}
\maketitle

\newpage
\tableofcontents

\newpage
\section{\textsf{2 - Time-Independent Schr\"{o}dinger Equation}}
\subsection{\textsf{2.1 - Stationary States}}
% QUESTION 1
\qs{2.1}{Prove the following three theorems:
\begin{enumerate}[label=(\alph*)]
	\item For normalizable solutions, the separation constant \( E \) must be \textit{real}. \textit{Hint}: Write \( E \) (in Equation 2.7) as \( E_{0} + i \Gamma \) (with \( E_{0} \) and \( \Gamma \) real), and show that if Equation 1.20 is to hold for all \( t \), \( \Gamma \) must be zero.
	\item The time-independent wave function \( \psi \!\left( x \right) \) can always be taken to be \textit{real} (unlike \( \Psi \!\left( x,t \right) \), which is necessarily complex). This doesn't mean that every solution to the time-independent Schr\"odinger equation \textit{is} real; what it says is that if you've got one that is \textit{not}, it can always be expressed as a linear combination of solutions (with the same energy) that \textit{are}. So you \textit{might as well} stick to \( \psi \)'s that are real. \textit{Hint}: If \( \psi \!\left( x \right) \) satisfies Equation 2.5, for a given \( E \), so too does its complex conjugate, and hence also the real linear combinations \( \!\left( \psi + \psi^{*} \right) \) and \( i \!\left( \psi - \psi^{*} \right) \).
	\item If \( V \!\left( x \right) \) is an \textbf{even function} (that is, \( V \!\left( -x \right) = V \!\left( x \right) \)) then \( \psi \!\left( x \right) \) can always be taken to be either even or odd. \textit{Hint}: If \( \psi \!\left( x \right) \) satisfies Equation 2.5, for a given \( E \), so too does \( \psi \!\left( -x \right) \), and hence also the even and odd linear combinations \( \psi \!\left( x \right) \pm \psi \!\left( -x \right) \).
\end{enumerate}}

\begin{enumerate}[label=(\alph*)]
	\item \begin{proof}[Proof]
	Let \( E = E_{0} + i \Gamma \), where \( E_{0}, \Gamma \in \mathbb{R} \). Then the wave function and its conjugate is
	\begin{align*}
		\Psi & = \psi \!\left( x \right)e^{-i \!\left( E_{0} + i \Gamma \right)t / \hbar} = \psi \!\left( x \right)e^{\!\left( \Gamma - iE_{0} \right)t / \hbar} \\
		\Psi^{*} & = \psi^{*} \!\left( x \right)e^{-i \!\left( -E_{0} + i \Gamma \right)t / \hbar} = \psi^{*} \!\left( x \right)e^{\!\left( \Gamma + iE_{0} \right)t / \hbar}
	\end{align*}
	Compute
	\begin{align*}
		\abs{\Psi \!\left( x,t \right)}^{2} & = \Psi^{*}\Psi = \psi^{*} \!\left( x \right)e^{\!\left( \Gamma + iE_{0} \right)t / \hbar} \cdot \psi \!\left( x \right)e^{\!\left( \Gamma - iE_{0} \right)t / \hbar} \\
						    & = \abs{\psi \!\left( x \right)}^{2}e^{2 \Gamma t / \hbar}
	\end{align*}
	By premise, \( \Psi \) is normalized. Then we must have
	\[
		\int^{\infty}_{-\infty} \abs{\Psi \!\left( x,t \right)}^{2} \, dx = \int^{\infty}_{-\infty} \abs{\psi \!\left( x \right)}^{2}e^{2 \Gamma t / \hbar} \, dx = e^{2 \Gamma t / \hbar}\int^{\infty}_{-\infty} \abs{\psi \!\left( x \right)}^{2} \, dx = 1
	\]
	Note that normalization (Equation 1.20) must hold for all \( t \). However, given that the exponential changes with \( t \) for non-zero \( \Gamma \), and the definite integral is a fixed non-zero value, we can conclude that we only have normalization when \( \Gamma = 0 \).
	\end{proof}
	\item \begin{proof}[Proof]
			Suppose that \( \psi \!\left( x \right) = a \!\left( x \right) + ib \!\left( x \right) \), where \( a \!\left( x \right), b \!\left( x \right) \in \mathbb{R} \), and that \( \psi \!\left( x \right) \) satisfies Equation 2.5:
	\begin{align*}
		-\frac{\hbar^{2}}{2m}\frac{d^2 \psi}{d x^2} + V \psi & = -\frac{\hbar^{2}}{2m}\frac{d^2 }{d x^2}\!\left[ a \!\left( x \right) + ib \!\left( x \right) \right] + V \!\left[ a \!\left( x \right) + ib \!\left( x \right) \right] \\
								     & = \!\left( -\frac{\hbar^{2}}{2m}\frac{d^2 a}{d x^2} + Va \right) + i\!\left( -\frac{\hbar^{2}}{2m}\frac{d^2 b}{d x^2} + Vb \right) \\
								     & = E \!\left[ a \!\left( x \right) + ib \!\left( x \right) \right] \\
								     & = E \psi
	\end{align*}
	Now, for the \textbf{same} \( E \) and \( V \) (which in the former case is real by part (a), and in the latter, assumed real throughout the text), \( a \!\left( x \right) \) and \( b \!\left( x \right) \) are both real solutions: \( V \) real guarantees the terms in the parentheses are real, and \( E \) real guarantees \( Ea \) and \( Eb \) are real. Then 
	\[
		-\frac{\hbar^{2}}{2m}\frac{d^2 a}{d x^2} + Va = Ea 
	\]
	and similarly for \( b \!\left( x \right) \). Thus \( \psi \) is itself a linear combination of two real solutions with the same energy \( E \).
	\end{proof}
	\item \begin{proof}[Proof]
	By premise, \( V \!\left( x \right) \) is even. Suppose \( \psi \) satisfies Equation 2.5 for a given \( E \). Then we can demonstrate that \( \psi \!\left( -x \right) \) also satisfies that equation. Define \( \phi \!\left( x \right) \equiv \psi \!\left( -x \right) \), noting that \( \phi''\!\left( x \right) = \psi''\!\left( -x \right)\!\left( -1 \right)^{2} = \psi''\!\left( -x \right) \) by the chain rule. Derive
	\begin{align*}
		-\frac{\hbar^{2}}{2m}\frac{d^2 \phi \!\left( x \right)}{d x^2} + V \!\left( x \right) \phi \!\left( x \right) & = -\frac{\hbar^{2}}{2m}\frac{d^2 \psi \!\left( -x \right)}{d x^2} + V \!\left( -x \right)\psi \!\left( -x \right) \\
															      & = E \psi \!\left( -x \right) = E \phi \!\left( x \right)
	\end{align*}
	where we invoke the premise that \( V \!\left( x \right) \) is even in the first equality. We now have Equation 2.5 at the point \( -x \). 

        Define \( \psi_{+}\!\left( x \right) = \psi \!\left( x \right) + \psi \!\left( -x \right) \) and \( \psi_{-}\!\left( x \right) = \psi \!\left( x \right) - \psi \!\left( -x \right) \). We can see that \( \psi_{+}\!\left( -x \right) = \psi \!\left( -x \right) + \psi \!\left( x \right) = \psi \!\left( x \right) + \psi \!\left( -x \right) = \psi_{+} \!\left( x \right) \), demonstrating that \( \psi_{+} \) is even. Analogously, \( -\psi_{-}\!\left( -x \right) = -\psi \!\left( -x \right) + \psi \!\left( x \right) = \psi_{-}\!\left( x \right) \), proving that \( \psi_{-}\!\left( x \right) \) is odd. 

	By linearity, we can take the linear combinations \( \psi \!\left( x \right) \pm \psi \!\left( -x \right) \), which themselves are solutions. Reconstructing \( \psi = \!\left( \psi_{+} + \psi_{-} \right) / 2 \), we establish that the even and odd components of every solution \( \psi \) of Equation 2.5, given that \( V \!\left( x \right) \) is even, are themselves solutions with the same energy \( E \).
	\end{proof}
\end{enumerate}

\newpage
% QUESTION 2
\qs{2.2}{Show that \( E \) must exceed the minimum value of \( V \!\left( x \right) \), for every normalizable solution to the time-independent Schr\"odinger equation. What is the classical analog to this statement? \textit{Hint}: Rewrite Equation 2.5 in the form
\[
	\frac{d^2 \psi}{d x^2} = \frac{2m}{\hbar^{2}}\!\left[ V \!\left( x \right) - E \right]\psi 
\]
if \( E < V_{\text{min}} \), then \( \psi \) and its second derivative always have the \textit{same sign} -- argue that such a function cannot be normalized.}
\begin{proof}[Proof]
	The geometric intuition is that when \( \psi \) and \( \psi'' \) share the same sign, the sign of \( \psi \) determines its curvature. In particular, when \( \psi > 0 \), it curves upwards in its neighborhood. The goal will be to prove that \( \psi \) is monotonic in two cases moving left- or rightwards, establishing divergence of \( \int \abs{\psi}^{2} \, dx \).

Suppose to the contrary that \( E \le V_{\text{min}} \). Then \( 2m \!\left[ V \!\left( x \right) - E \right] / \hbar^{2} \ge 0 \), and \( \psi''\psi \ge 0 \). We must now deduce the behavior of \( \psi \) from this fact. By problem 2.1, we may take \( \psi \) to be real. Fix a point \( x_{0} \), assume without loss of generality \( \psi \!\left( x_{0} \right) > 0 \) (from which it immediately follows \( \psi''\!\left( x_{0} \right) \ge 0 \)), and consider two cases where \( \psi'\!\left( x_{0} \right) \le 0 \) and \( \psi'\!\left( x_{0} \right) > 0 \).

\textbf{Case 1:} \( \psi'\!\left( x_{0} \right) \le 0 \). With \( x_{0} \ge x \rightarrow -\infty \), since we must have \( \psi''\!\left( x \right) \ge 0 \) (to the left of \( x_{0} \)), we first determine the behavior  of \( \psi \!\left( x \right) \). Consider the nearest point \( x_{1} < x_{0} \) such that \( \psi \!\left( x_{1} \right) \le 0 < \psi \!\left( x_{0} \right) \). By the Mean Value Theorem, we must have \( \psi' > 0 \) at some point \( c \in (x_{1},x_{0}] \). But since we selected \( x_{1} \) such that it is the \textit{first} point to the left of \( x_{0} \) such that \( \psi \!\left( x_{1} \right) \le 0 \), we can now claim \( \psi > 0 \), and thus \( \psi'' \ge 0 \), on \( (x_{1},x_{0}] \). However, because \( \psi'\!\left( c \right) > 0 \) and \( \psi'\!\left( x_{0} \right) \le 0 \), a second application of the Mean Value Theorem reveals that \( \psi'' < 0 \) \textit{somewhere} on \( (x_{1},x_{0}] \), a contradiction. Then it must be that \( \psi > 0 \) and \( \psi'' \ge 0 \) for \textit{all} points \( \!\left( -\infty,x_{0} \right) \), with \( \psi' \le \psi'\!\left( x_{0} \right) \le 0 \) -- monotonicity. Simply: \( \psi \) is always curving away from the horizontal axis on that interval, and with monotonicity, it follows that 
\[ 
	\int^{x_{0}}_{-\infty} \abs{\psi \!\left( x \right)}^{2} \, dx \ge \abs{\psi \!\left( x_{0} \right)}^{2}\int^{x_{0}}_{-\infty} \, dx  = \infty 
\]

\textbf{Case 2:} \( \psi'\!\left( x_{0} \right) > 0 \). Reverse the argument in case 1 going rightward to conclude \( \psi > 0 \), \( \psi'' \ge 0 \), and \( \psi' \ge \psi'\!\left( x_{0} \right) > 0 \). Similarly, with monotonicity:
\[
	\int^{\infty}_{x_{0}} \abs{\psi \!\left( x \right)}^{2} \, dx \ge \abs{\psi \!\left( x_{0} \right)}^{2}\int^{\infty}_{x_{0}}  \, dx = \infty  
\]

In either case, \( \psi \) cannot be normalized, a contradiction.

\begin{center}
\begin{tikzpicture}[
	>=Stealth,
	axis/.style={->, thick},
	curve/.style={deepblue, very thick},
	tangent/.style={deepred, thick},
	marker/.style={dashed, gray},
	region/.style={fill=deepblue!10},
	every node/.style={font=\small}
]
	% Case 1: psi(x0) > 0, psi'(x0) <= 0, convex everywhere, so psi blows up as x -> -infinity
	\begin{scope}[scale=0.8]
		\coordinate (X0) at (1,0);
		\coordinate (P0) at (1,0.5744);
		% divergent area under |psi|^2 on (-infty, x0]
		\fill[region] plot[domain=-2.7:1, samples=60] (\x,{0.5*exp(-0.6*\x)+0.3}) -- (X0) -- (-2.7,0) -- cycle;
		% axes
		\draw[axis] (-3.2,0) -- (3.2,0) node[below] {$x$};
		\draw[axis] (0,-0.3) -- (0,3.3) node[left] {$\psi$};
		% the solution
		\draw[curve, domain=-2.7:3, samples=80] plot (\x,{0.5*exp(-0.6*\x)+0.3});
		\node[right] at (-1.9,2.1) {$\psi\!\left( x \right)$};
		\node[right] at (-2.5,3.0) {$\psi \rightarrow \infty$};
		\node at (1.8,1.4) {$\psi'' > 0$};
		% the point x0 and the tangent there
		\draw[marker] (X0) -- (P0);
		\draw[tangent, domain=0:2.2] plot (\x,{0.5744-0.1646*(\x-1)});
		\node[tangent] at (2.3,0.8) {$\psi'\!\left( x_{0} \right) \le 0$};
		\fill (P0) circle (1.6pt);
		\node[below] at (X0) {$x_{0}$};
		\node[below] at (-1.4,-0.3) {$\int_{-\infty}^{x_{0}} \abs{\psi}^{2} \, dx = \infty$};
		\node[above] at (0,3.5) {\textbf{Case 1:} $\psi'\!\left( x_{0} \right) \le 0$};
	\end{scope}

	% Case 2: psi'(x0) > 0, so psi blows up as x -> +infinity
	\begin{scope}[xshift=7cm, scale=0.8]
		\coordinate (X0) at (-1,0);
		\coordinate (P0) at (-1,0.5744);
		% divergent area under |psi|^2 on [x0, infty)
		\fill[region] (X0) -- plot[domain=-1:2.7, samples=60] (\x,{0.5*exp(0.6*\x)+0.3}) -- (2.7,0) -- cycle;
		% axes
		\draw[axis] (-3.2,0) -- (3.2,0) node[below] {$x$};
		\draw[axis] (0,-0.3) -- (0,3.3) node[left] {$\psi$};
		% the solution
		\draw[curve, domain=-3:2.7, samples=80] plot (\x,{0.5*exp(0.6*\x)+0.3});
		\node[left] at (1.9,2.1) {$\psi\!\left( x \right)$};
		\node[left] at (2.5,3.0) {$\psi \rightarrow \infty$};
		\node at (-1.8,1.4) {$\psi'' > 0$};
		% the point x0 and the tangent there
		\draw[marker] (X0) -- (P0);
		\draw[tangent, domain=-2.2:0] plot (\x,{0.5744+0.1646*(\x+1)});
		\node[tangent] at (-2.3,0.8) {$\psi'\!\left( x_{0} \right) > 0$};
		\fill (P0) circle (1.6pt);
		\node[below] at (X0) {$x_{0}$};
		\node[below] at (1.4,-0.3) {$\int_{x_{0}}^{\infty} \abs{\psi}^{2} \, dx = \infty$};
		\node[above] at (0,3.5) {\textbf{Case 2:} $\psi'\!\left( x_{0} \right) > 0$};
	\end{scope}
\end{tikzpicture}
\end{center}
\end{proof}

In classical mechanics, the total energy of a particle \textbf{cannot} be less than the potential energy, but it may be equal in the case where the particle is at rest. The Hamiltonian is
\[
	H \!\left( x,p \right) = \underbrace{\frac{p^{2}}{2m}}_{\text{kinetic}} + \underbrace{V \!\left( x \right)}_{\text{potential}}
\]
in which the kinetic energy is either positive or zero. Thus in the classical case, \( E \ge V_{\text{min}} \), but in the quantum mechanical instance, \( E > V_{\text{min}} \). We call this \textbf{zero-point energy}.

\subsection{\textsf{2.2 - The Infinite Square Well}}
% QUESTION 3
\qs{2.3}{Show that there is no acceptable solution to the (time-independent) Schr\"odinger equation for the infinite square well with \( E = 0 \) or \( E < 0 \). (This is a special case of the general theorem in Problem 2.2, but this time do it by explicitly solving the Schr\"odinger equation, and showing that you cannot meet the boundary conditions.)}

\begin{proof}[Proof]
Consider the time-independent Schr\"odinger equation
\[
	-\frac{\hbar^{2}}{2m}\frac{d^2 \psi}{d x^2} = E \psi 
\]
Consider two cases subject to the boundary conditions \( \psi \!\left( 0 \right) = \psi \!\left( a \right) = 0 \).

\underline{\( E = 0 \)}: The ODE simplifies to
\[
	\frac{d^2 \psi}{d x^2} = 0 
\]
where integrating twice produces solution
\[
	\psi \!\left( x \right) = c_{1}x + c_{2} 
\]
Imposing the boundary conditions forces \( \psi \!\left( x \right) = 0 \), which is ineligible for non-normalizability:
\begin{align*}
	\psi \!\left( 0 \right) & = c_{2} = 0 \\
	\psi \!\left( a \right) & = c_{1}a = 0 \quad \implies \quad c_{1} = 0
\end{align*}

\underline{\( E < 0 \)}: Rewrite the ODE as 
\[
	\frac{d^2 \psi}{d x^2} = \lambda \psi, \quad \lambda = -\frac{2mE}{\hbar^{2}} 
\]
where it is apparent that \( \lambda > 0 \), giving solution
\[
	\psi \!\left( x \right) = c_{1}\cosh\!\left(\sqrt{\lambda}x\right) + c_{2}\sinh\!\left(\sqrt{\lambda}x\right) 
\]
The boundary conditions require
\begin{align*}
	\psi \!\left( 0 \right) & = c_{1} = 0 \\
	\psi \!\left( a \right) & = c_{2}\sinh\!\left(\sqrt{\lambda}a\right) = 0 \quad \implies \quad c_{2} = 0
\end{align*}
where the implication in the second condition arises from the fact that \( \sqrt{\lambda}a \neq 0 \), and \( \sinh\!\left(x\right) = 0 \) only when \( x = 0 \), thus our only option is to have \( c_{2} = 0 \). Once more, the trivial -- and non-normalizable -- solution \( \psi \!\left( x \right) = 0 \) remains.

Ergo, it is impossible to have \( E = 0 \) or \( E < 0 \).
\end{proof}

% QUESTION 4
\qs{2.4}{Calculate \( \braket{x} \), \( \braket{x^{2}} \), \( \braket{p} \), \( \braket{p^{2}} \), \( \sigma_{x} \), and \( \sigma_{p} \), for the \( n \)th stationary state of the infinite square well. Check that the uncertainty principle is satisfied. Which state comes closest to the uncertainty limit?}

The \( n \)th stationary state (and its conjugate) of the infinite square well is
\[
	\psi_{n}\!\left( x \right) = \psi_{n}\!\left( x \right)^{*} = \sqrt{\frac{2}{a}}\sin\!\left(\frac{n \pi}{a}x\right) 
\]
To make our lives easier when computing integrals with integrands \( x^{2}\sin\!\left(kx\right) \) or \( x^{2}\cos\!\left(kx\right) \), use the tricks
\begin{align*}
	\int^{a}_{0} x^{2}\sin\!\left(kx\right) \, dx & = -\frac{d^2 }{d k^2}\int^{a}_{0} \sin\!\left(kx\right) \, dx = \frac{d^2 }{d k^2}\!\left[ \frac{\cos\!\left(ka\right) - 1}{k} \right] \\
	\int^{a}_{0} x^{2}\cos\!\left(kx\right) \, dx & = -\frac{d^2 }{d k^2}\int^{a}_{0} \cos\!\left(kx\right) \, dx = -\frac{d^2 }{d k^2}\!\left[ \frac{\sin\!\left(ka\right)}{k} \right]  
\end{align*}
which conveniently eliminates the need for integration by parts! Now compute 
\begin{align*}
\braket{x} & = \frac{2}{a}\int^{a}_{0} x \sin^{2}\!\left(\frac{n \pi}{a}x\right) \, dx  = \frac{a}{2} \quad (\text{by symmetry of } \abs{\psi_{n}}^{2} \text{ about } a/2) \\
	\braket{x^{2}} & = \frac{2}{a}\int^{a}_{0} x^{2}\sin^{2}\!\left(\frac{n \pi}{a}x\right) \, dx \\
		       & = \frac{1}{a}\!\left[ \int^{a}_{0} x^{2} \, dx - \int^{a}_{0} x^{2}\cos\!\left(\frac{2n \pi}{a}x\right) \, dx \right] \\
		       & = \frac{a^{2}}{3} + \frac{1}{a}\frac{d^2 }{d k^2}\!\left[ \frac{\sin\!\left(ka\right)}{k} \right] \quad (k = 2n \pi / a) \\
		       & = \frac{a^{2}}{3} - \frac{a\sin\!\left(ka\right)}{k} - \frac{2 \cos\!\left(ka\right)}{k^{2}} + \frac{2 \sin\!\left(ka\right)}{ak^{3}} \\
		       & = \frac{a^{2}}{3} - \frac{a^{2}\sin\!\left(2n \pi\right)}{2n \pi} - \frac{a^{2} \cos\!\left(2n \pi\right)}{2n^{2}\pi^{2}} + \frac{2a^{2} \sin\!\left(2n \pi\right)}{8n^{3}\pi^{3}} \\
		       & = a^{2}\!\left( \frac{1}{3} - \frac{1}{2n^{2}\pi^{2}} \right) \\
	    \braket{p} & = \frac{2}{a}\int^{a}_{0} \sin\!\left(\frac{n \pi}{a}x\right)\frac{\hbar}{i}\frac{d }{d x}\sin\!\left(\frac{n \pi}{a}x\right) \, dx \quad (\text{"cancel" the } dx \text{'s and treat } d \sin\!\left(n \pi x / a\right) \text{ as a differential}) \\
		       & = \frac{\hbar}{ai}\sin^{2}\!\left(\frac{n \pi}{a}x\right)\bigg|^{a}_{0} = 0 \\
	\braket{p^{2}} & = \frac{2}{a}\int^{a}_{0} \sin\!\left(\frac{n \pi}{a}x\right)\!\left( \frac{\hbar}{i} \right)^{2}\frac{d^2 }{d x^2}\sin\!\left(\frac{n \pi}{a}x\right) \, dx \\
		       & = \frac{2n^{2}\pi^{2} \hbar^{2}}{a^{3}}\int^{a}_{0} \sin^{2}\!\left(\frac{n \pi}{a}x\right) \, dx = \!\left( \frac{n \pi \hbar}{a} \right)^{2} \\
	\sigma_{x} & = \sqrt{\braket{x^{2}} - \braket{x}^{2}} = a \sqrt{\frac{1}{12} - \frac{1}{2n^{2}\pi^{2}}} = \frac{a}{2n \pi}\sqrt{\frac{n^{2}\pi^{2} - 6}{3}} \\
	    \sigma_{p} & = \sqrt{\braket{p^{2}} - \braket{p}^{2}} = \frac{n \pi \hbar}{a}
\end{align*}

Multiply the standard deviations to check
\[
	\sigma_{x}\sigma_{p} = \frac{\hbar}{2}\sqrt{\frac{n^{2}\pi^{2} - 6}{3}} > \frac{\hbar}{2}
\]
where the state \( n = 1 \) comes closest to the uncertainty limit, namely 
\[
	\frac{\hbar}{2}\sqrt{\frac{\pi^{2} - 6}{3}} \approx 1.14 \cdot \frac{\hbar}{2} 
\]

% QUESTION 2.5
\qs{2.5}{A particle in the infinite square well has as its initial wave function an even mixture of the first two stationary states:
\[
	\Psi \!\left( x,0 \right) = A \!\left[ \psi_{1}\!\left( x \right) + \psi_{2}\!\left( x \right) \right]
\]
\begin{enumerate}[label=(\alph*)]
	\item Normalize \( \Psi \!\left( x,0 \right) \). (That is, find \( A \). This is very easy, if you exploit the orthonormality of \( \psi_{1} \) and \( \psi_{2} \). Recall that, having normalized \( \Psi \) at \( t = 0 \), you can rest assured that it \textit{stays} normalized--if you doubt this, check it explicitly after doing part (b).)
	\item Find \( \Psi \!\left( x,t \right) \) and \( \abs{\Psi \!\left( x,t \right)}^{2} \). Express the latter as a sinusoidal function of time, as in Example 2.1. To simplify the result, let \( \omega \equiv \pi^{2}\hbar / 2ma^{2} \).
	\item Compute \( \braket{x} \). Notice that it oscillates in time. What is the angular frequency of the oscillation? What is the amplitude of the oscillation? (If your amplitude is greater than \( a / 2 \), go directly to jail.)
	\item Compute \( \braket{p} \). (As Peter Lorre would say, ``Do it ze \textit{kveek} vay, Johnny!'')
	\item If you measured the energy of this particle, what values might you get, and what is the probability of getting each of them? Find the expectation value of \( H \). How does it compare with \( E_{1} \) and \( E_{2} \)?
\end{enumerate}}

\begin{enumerate}[label=(\alph*)]
	\item Compute
	\[
		\abs{\Psi \!\left( x,0 \right)}^{2} = \abs{A}^{2}\!\left[ \psi^{2}_{1} + 2 \psi_{1}\psi_{2} + \psi^{2}_{2} \right]
	\]
	where \( \psi_{1} \) and \( \psi_{2} \) are real, allowing us to drop the conjugates as \( \psi^{*} = \psi \). Then we may derive (using orthonormality)
	\begin{align*}
		\int^{\infty}_{-\infty} \abs{\Psi \!\left( x,0 \right)}^{2} \, dx & = \abs{A}^{2}\!\left[ \underbrace{\int^{\infty}_{-\infty} \psi^{2}_{1} \, dx}_{= 1} + 2 \underbrace{\int^{\infty}_{-\infty} \psi_{1}\psi_{2} \, dx}_{= 0} + \underbrace{\int^{\infty}_{-\infty} \psi^{2}_{2} \, dx}_{= 1} \right] \\
		 & = 2\abs{A}^{2} = 1
	\end{align*}
	Implying \( \abs{A} = 1 / \sqrt{2} \); choose the positive branch to conclude \( A = 1 / \sqrt{2} \).
	\item Adding in the time dependence, we have
	\[
		\Psi \!\left( x,t \right) = \frac{1}{\sqrt{2}}\!\left[ \psi_{1}e^{-iE_{1}t / \hbar} + \psi_{2}e^{-iE_{2}t / \hbar} \right] 
	\]
	where \( E_{n} = n^{2}\pi^{2}\hbar^{2} / 2ma^{2} \). Then
	\[
		\Psi^{*}\!\left( x,t \right) = \frac{1}{\sqrt{2}}\!\left[ \psi_{1}e^{iE_{1}t / \hbar} + \psi_{2}e^{iE_{2}t / \hbar} \right]
	\]
	enabling us to compute
	\begin{align*}
		\abs{\Psi \!\left( x,t \right)}^{2} & = \Psi^{*}\Psi \\
						    & = \frac{1}{2}\!\left[ \psi_{1}e^{iE_{1}t / \hbar} + \psi_{2}e^{iE_{2}t / \hbar} \right]\!\left[ \psi_{1}e^{-iE_{1}t / \hbar} + \psi_{2}e^{-iE_{2}t / \hbar} \right] \\ 
						    & = \frac{1}{2}\!\left[ \psi^{2}_{1} + \psi^{2}_{2} + \psi_{1}\psi_{2}\!\left( e^{i \!\left( E_{2} - E_{1} \right)t / \hbar} + e^{-i \!\left( E_{2} - E_{1} \right)t / \hbar} \right) \right] \\
						    & = \frac{1}{2}\!\left[ \psi^{2}_{1} + \psi^{2}_{2} + 2\psi_{1}\psi_{2}\cos\!\left(\frac{\!\left( E_{2} - E_{1} \right)t}{\hbar}\right) \right] \\
						    & = \frac{1}{2}\!\left[ \psi^{2}_{1} + \psi^{2}_{2} + 2 \psi_{1}\psi_{2}\cos\!\left(\frac{3 \pi^{2}\hbar}{2ma^{2}}t\right) \right] \\
						    & = \frac{1}{2}\!\left( \psi^{2}_{1} + \psi^{2}_{2} \right) + \psi_{1}\psi_{2}\cos\!\left(3 \omega t\right)
	\end{align*}
	where \( \omega = \pi^{2}\hbar / 2ma^{2} \).
	\item Compute
	\begin{align*}
		\braket{x}\!\left( t \right) & = \int^{\infty}_{-\infty} x \abs{\Psi \!\left( x,t \right)}^{2} \, dx \\
			   & = \frac{1}{2}\!\left[ \int^{\infty}_{-\infty} x \psi^{2}_{1} \, dx+ \int^{\infty}_{-\infty} x \psi^{2}_{2} \, dx \right] + \cos\!\left(3 \omega t\right)\int^{\infty}_{-\infty} x \psi_{1}\psi_{2} \, dx \\ 
			   & = \frac{a}{2} + \frac{2}{a}\cos\!\left(3 \omega t\right)\int^{a}_{0} x \sin\!\left(\frac{\pi x}{a}\right)\sin\!\left(\frac{2 \pi x}{a}\right) \, dx
	\end{align*}
	Note that we change the integral bounds to \( \!\left[ 0,a \right] \) as the wave functions are explicitly defined on those bounds. Using the product-to-sum identity, continue with
	\[
		 = \frac{a}{2} + \frac{\cos\!\left(3 \omega t\right)}{a}\int^{a}_{0} x \!\left[ \cos\!\left(\frac{\pi x}{a}\right) - \cos\!\left(\frac{3 \pi x}{a}\right) \right] \, dx
	\]
	And then use the trick from problem 2.4 to see
	\begin{align*}
		\int^{a}_{0} x \cos\!\left(kx\right) \, dx = \frac{d }{d k}\int^{a}_{0} \sin\!\left(kx\right) \, dx & = -\frac{d }{d k}\!\left[ \frac{\cos\!\left(ka\right) - 1}{k} \right] \\
							   & = - \!\left[ \frac{1 - \cos\!\left(ka\right)}{k^{2}} - \frac{a \sin\!\left(ka\right)}{k} \right] \\
							   & = \frac{ak \sin\!\left(ka\right) + \cos\!\left(ka\right) - 1}{k^{2}}
	\end{align*}
	Evaluating for \( k = \pi / a \) and \( k = 3 \pi / a \), which causes the sines to vanish, compute
	\begin{align*}
		 & = \frac{a}{2} + \frac{\cos\!\left(3 \omega t\right)}{a}\!\left[ \!\left( \frac{\cos\!\left(\pi\right) - 1}{\pi^{2} / a^{2}} \right) - \!\left( \frac{\cos\!\left(3 \pi\right) - 1}{9 \pi^{2} / a^{2}} \right) \right] \\
		 & = \frac{a}{2} + \frac{\cos\!\left(3 \omega t\right)}{a}\!\left[ -\frac{2a^{2}}{\pi^{2}} + \frac{2a^{2}}{9 \pi^{2}} \right] \\ 
		 & = \frac{a}{2} - \frac{16a}{9 \pi^{2}}\cos\!\left(3 \omega t\right)
	\end{align*}
	Since \( 16 / 9 \pi^{2} \approx 0.18 \), our amplitude \( 16a / 9 \pi^{2} < a/2 \). The angular frequency is \( 3 \omega \). Now, evaluate \( \braket{x} \) at \( t = 0 \):
	\[
		\braket{x}\!\left( 0 \right) = 
	\]
	Observe further that \( \psi_{1}, \psi_{2} > 0 \) on \( \!\left[ 0,a/2 \right] \), while \( \psi_{1} > 0 \) but \( \psi_{2} < 0 \) on \( \!\left[ a/2,a \right] \). It follows that the density \( \abs{\Psi \!\left( x,t \right)}^{2} \) is heaviest along \( \!\left[ 0,a/2 \right] \), 
	\item Compute, using the result that \( \braket{p} = m d \braket{x} / dt \):
	\begin{align*}
		\braket{p} & = m \frac{d \braket{x}}{d t} \\
			   & = \frac{16ma}{3 \pi^{2}} \omega \sin\!\left(3 \omega t\right) \\
			   & = \frac{16ma}{3 \pi^{2}}\frac{\pi^{2}\hbar}{2ma^{2}}\sin\!\left(3 \omega t\right) \\
			   & = \frac{8 \hbar}{3a}\sin\!\left(3 \omega t\right)
	\end{align*}
	\item We compute the probabilities as the squared Fourier coefficients of projecting \( \psi_{n} \) onto \( \Psi \!\left( x,0 \right) \). Compute for \( n = 1, 2 \)
	\[
		c_{n} = \int^{a}_{0} \psi_{n}\Psi \!\left( x,0 \right) \, dx = A\int^{a}_{0} \psi_{n}\!\left[ \psi_{1} + \psi_{2} \right] \, dx = A \underbrace{\int^{a}_{0} \psi^{2}_{n} \, dx}_{= 1} = A
	\]
	Concretely, these coefficients are
	\[
		c_{1} = c_{2} = \frac{1}{\sqrt{2}}
	\]
	from which we can confirm
	\[
		\abs{c_{1}}^{2} + \abs{c_{2}}^{2} = \!\left( \frac{1}{\sqrt{2}} \right)^{2} + \!\left( \frac{1}{\sqrt{2}} \right)^{2} = 1 
	\]
	Accordingly, we can now conclude that we have a \( 1 / 2 \) probability of measuring the particle energy as \( E_{1} \) or \( E_{2} \). The expectation of the energy \( H \) is 
	\begin{align*}
		\braket{H} & = \abs{c_{1}}^{2}E_{1} + \abs{c_{2}}^{2}E_{2} \\
			   & = \frac{1}{2} \cdot \frac{\pi^{2}\hbar^{2}}{2ma^{2}} + \frac{1}{2} \cdot \frac{4 \pi^{2}\hbar^{2}}{2ma^{2}} \\
			   & = \frac{5 \pi^{2}\hbar^{2}}{4ma^{2}}
	\end{align*}
	which is the midpoint value in between \( E_{1} \) and \( E_{2} \), as anticipated as the probabilities of either state are equal.
\end{enumerate}

% QUESTION 6
\qs{2.6}{Although the \textit{overall} phase constant of the wave function is of no physical significance (it cancels out whenever you calculate a measurable quantity), the \textit{relative} phase of the coefficients in Equation 2.17 \textit{does} matter. For example, suppose we change the relative phase of \( \psi_{1} \) and \( \psi_{2} \) in Problem 2.5:
\[
	\Psi \!\left( x,0 \right) = A \!\left[ \psi_{1}\!\left( x \right) + e^{i \phi}\psi_{2}\!\left( x \right) \right] 
\]
where \( \phi \) is some constant. Find \( \Psi \!\left( x,t \right) \), \( \abs{\Psi \!\left( x,t \right)}^{2} \), and \( \braket{x} \), and compare your results with what you got before. Study the special cases \( \phi = \pi / 2 \) and \( \phi = \pi \). (For a graphical exploration of this problem see the applet in footnote 7.)}

Using the wave function conjugate


\end{document}
